\documentclass{article}
\usepackage{amsmath, amssymb, amsthm}

\newtheorem{theorem}{Theorem}[section]
\newtheorem{lemma}[theorem]{Lemma}
\newtheorem{corollary}[theorem]{Corollary}
\theoremstyle{definition}
\newtheorem{definition}[theorem]{Definition}

\title{Irrationality of $\sqrt{2}$ and Friends}
\author{Fermat Demo}
\date{}

\begin{document}
\maketitle

% ── Open this file in Fermat and click "Submit All" to see the full pipeline.
% ── The Hard theorem will go through sketch → prove → verify before review.

\section{Setup}

\begin{definition}[Rational Number]
\label{def:rational}
A real number $x$ is \emph{rational} if there exist integers $p, q$ with
$q \neq 0$ such that $x = p/q$.  We say $x$ is \emph{irrational} otherwise.
\end{definition}

\begin{lemma}[Parity lemma]
\label{lem:parity}
For any integer $n$, if $n^2$ is even then $n$ is even.
\end{lemma}
% [PROVE IT: Easy]

\section{Main Results}

\begin{theorem}[Irrationality of $\sqrt{2}$]
\label{thm:sqrt2}
$\sqrt{2}$ is irrational.
\end{theorem}
% [PROVE IT: Easy]
% SKETCH: Assume sqrt(2) = p/q in lowest terms. Then 2q^2 = p^2, so p^2 is
%   even, hence p is even by Lemma~\ref{lem:parity}. Write p = 2k; then
%   2q^2 = 4k^2, so q^2 = 2k^2, meaning q is also even — contradicting lowest terms.

\begin{theorem}[Irrationality of $\sqrt{p}$ for primes $p$]
\label{thm:sqrt-prime}
Let $p$ be a prime number.  Then $\sqrt{p}$ is irrational.
\end{theorem}
% [PROVE IT: Medium]

\begin{corollary}[Density of irrationals]
\label{cor:dense}
Between any two distinct real numbers there exists an irrational number.
\end{corollary}
% [PROVE IT: Hard]
% SKETCH: Let a < b. WLOG assume a, b >= 0 (the negative case is symmetric).
%   The interval (a/sqrt(2), b/sqrt(2)) contains a rational r = p/q (by density
%   of rationals). Then r*sqrt(2) = p*sqrt(2)/q lies in (a, b) and is irrational
%   by Theorem~\ref{thm:sqrt2} (since p != 0).

\end{document}
